\chapter{ماژول‌ها و بسته‌ها}%
\label{ch:modules}

تا اینجا همهٔ برنامه‌ها را در یک فایل نوشتیم. اما برنامه‌های واقعی بزرگ‌اند و
نمی‌توان همه را در یک فایل گنجاند. \textbf{ماژول‌ها} و \textbf{بسته‌ها} به ما
اجازه می‌دهند برنامه را به فایل‌های کوچک‌تر و معنادار بشکنیم و کدها را میانِ
آن‌ها به اشتراک بگذاریم.

\section{چرا برنامه را بشکنیم؟}

تصور کنید یک برنامهٔ بزرگ دارید: بخشی برای محاسباتِ ریاضی، بخشی برای کار با
متن، بخشی برای رابطِ کاربری. اگر همه را در یک فایل بریزید، یافتنِ هر چیز دشوار
می‌شود. با تقسیمِ برنامه به ماژول‌ها:
\begin{itemize}
  \item هر فایل یک مسئولیتِ روشن دارد.
  \item کدها را می‌توان میانِ پروژه‌ها به اشتراک گذاشت و دوباره به‌کار برد.
  \item چند نفر می‌توانند هم‌زمان روی بخش‌های مختلف کار کنند.
\end{itemize}

\section{فراخوانیِ کتابخانه‌های آماده}

ساده‌ترین شکلِ ماژول، کتابخانه‌های آمادهٔ خودِ سلام‌اند. با واژهٔ
\code{واردسازی} (معادلِ \code{import}) آن‌ها را به برنامه می‌آوریم. این مثالِ
واقعی ماژولِ ریاضی را وارد می‌کند:

\begin{salamfa}
واردسازی "ریاضی"

روال ریشه:
    سرچاپ ریاضی.توان(2.0, 10.0)
    سرچاپ ریاضی.جذر(144.0)
پایان
\end{salamfa}

پس از \code{واردسازی "ریاضی"}، می‌توانیم با نامِ ماژول و یک نقطه، به روال‌هایش
دسترسی پیدا کنیم: \code{ریاضی.توان(...)}، \code{ریاضی.جذر(...)}. این نقطه می‌گوید
«توانِ متعلق به ماژولِ ریاضی».

\begin{note}
نامِ ماژول‌های استاندارد هم مانندِ همه‌چیزِ دیگر در سلام دوزبانه است.
\code{واردسازی "ریاضی"} و \code{واردسازی "math"} یک کار می‌کنند، و می‌توانید با
\code{ریاضی.جذر} یا \code{math.Sqrt} به روال‌هایش برسید. در فصلِ ۱۵ کتابخانهٔ
استاندارد را به‌تفصیل می‌بینیم.
\end{note}

\section{ساختنِ ماژولِ خودمان}

می‌توانیم کدهای خود را در یک فایلِ جدا بگذاریم و در فایلِ دیگر از آن‌ها استفاده
کنیم. نخست یک فایلِ کتابخانه می‌سازیم. این مثالِ واقعی، فایلی به نامِ
\code{\_falib.salam} است:

\begin{salamfa}
بسته falib
همگانی روال جمع(a : صحیح, b : صحیح) : صحیح: برگشت a + b پایان
همگانی پایا پاسخ = 42
\end{salamfa}

دو نکتهٔ کلیدی:
\begin{itemize}
  \item \code{بسته falib} (معادلِ \code{package}) به این فایل یک \emph{نام}
        می‌دهد و می‌گوید که این یک بستهٔ کتابخانه‌ای است، نه یک برنامهٔ اجرایی.
  \item \code{همگانی} (معادلِ \code{pub}) پیش از \code{روال جمع} و \code{پایا
        پاسخ} می‌گوید که این‌ها \emph{بیرون} از فایل هم در دسترس‌اند.
\end{itemize}

\subsection{همگانی در برابرِ خصوصی}

به‌طورِ پیش‌فرض، هر چیزی که در یک ماژول تعریف می‌کنید \emph{خصوصی} است؛ یعنی
فقط درونِ همان فایل دیده می‌شود. تنها چیزهایی که با \code{همگانی} نشانه‌گذاری
کنید، از بیرون قابلِ دسترس‌اند. این یک اصلِ مهمِ طراحی است: \emph{فقط آنچه را
که دیگران واقعاً لازم دارند، همگانی کنید} و باقی را خصوصی نگه دارید.

\begin{tip}
چرا چیزها را خصوصی نگه داریم؟ چون هرچه کمتر همگانی باشد، آزادیِ شما برای تغییرِ
درونِ ماژول بیشتر است. اگر همه‌چیز همگانی باشد، هر تغییرِ کوچک ممکن است کدِ
دیگران را بشکند. مرزِ روشنِ «همگانی/خصوصی» به شما اجازه می‌دهد بی‌دغدغه درونِ
ماژول را بهبود دهید.
\end{tip}

\subsection{استفاده از ماژولِ خودمان}

حالا در فایلِ اصلی، ماژول را وارد و از آن استفاده می‌کنیم. این مثالِ واقعی نشان
می‌دهد چگونه با یک \emph{نامِ مستعار} ماژول را وارد کنیم:

\begin{salamfa}
واردسازی fl "_falib.salam"

روال ریشه:
    سرچاپ fl.جمع(2, 3)
    سرچاپ fl.پاسخ
پایان
\end{salamfa}

اینجا \code{واردسازی fl "\_falib.salam"} فایل را وارد می‌کند و به آن نامِ
کوتاهِ \code{fl} می‌دهد. سپس با \code{fl.جمع(2, 3)} و \code{fl.پاسخ} به
عضوهای همگانیِ آن دسترسی پیدا می‌کنیم. خروجی: \code{5} و \code{42}.

\begin{note}
نامِ مستعار (اینجا \code{fl}) به‌ویژه وقتی مفید است که نامِ ماژول طولانی باشد یا
دو ماژول نامِ یکسان داشته باشند. می‌توانید هر نامِ کوتاه و گویایی برگزینید.
\end{note}

\section{دوزبانه‌سازیِ ماژول‌های خودمان}

یکی از زیباترین ویژگی‌های سلام این است که می‌توانید برای کدِ \emph{خودتان} هم
نام‌های فارسی و انگلیسی تعریف کنید. این کار با حاشیه‌نویسی‌های \code{@en} و
\code{@fa} انجام می‌شود. این مثالِ واقعی یک ساختارِ کاملاً دوزبانه می‌سازد:

\begin{salamfa}
@en "Point"
@fa "نقطه"
struct Point:
    @en "x"
    @fa "ایکس"
    x : int = 0
    @en "y"
    @fa "وای"
    y : int = 0
    @en "Sum"
    @fa "جمع"
    func Sum() : int: ret this.x + this.y end
end
روال ریشه:
    mut p Point = Point { ایکس = 3, وای = 4 }
    چاپ : p.ایکس
    چاپ : p.جمع()
    چاپ : lang()
پایان:
\end{salamfa}

با \code{@en "Point"} و \code{@fa "نقطه"} به ساختار دو نام دادیم. حالا در یک
برنامهٔ فارسی می‌توان میدان‌ها و متدها را با نامِ فارسی صدا زد
(\code{p.ایکس}، \code{p.جمع()})، هرچند خودِ ساختار با نام‌های انگلیسی تعریف شده
است. روالِ ویژهٔ \code{lang()} نیز زبانِ جاری را برمی‌گرداند (\code{"fa"} یا
\code{"en"}).

\begin{deep}[نگاهی به درون: چگونه یک کد، دو زبان دارد؟]
وقتی شما \code{@fa "نقطه"} می‌نویسید، کامپایلر یک \emph{نگاشت} میانِ نامِ فارسی
و نامِ انگلیسیِ اصلی می‌سازد. در زمانِ کامپایل، هر جا که \code{نقطه} ببیند، آن
را به همان \code{Point}ِ اصلی ترجمه می‌کند. پس در کدِ نهایی تنها یک گونه وجود
دارد؛ دو نام فقط دو \emph{درگاهِ ورودی} به همان چیزند. به همین دلیل است که
دوزبانه‌بودن هیچ هزینهٔ اجرایی ندارد همه‌چیز در زمانِ کامپایل حل می‌شود.
\end{deep}

\begin{tip}
این ویژگی برای نوشتنِ کتابخانه‌هایی که هم فارسی‌زبانان و هم انگلیسی‌زبانان از
آن‌ها استفاده می‌کنند فوق‌العاده است. یک‌بار کتابخانه را با \code{@en} و
\code{@fa} می‌نویسید و هر کاربر با زبانِ خودش از آن بهره می‌برد درست همان
کاری که خودِ کتابخانهٔ استانداردِ سلام انجام می‌دهد.
\end{tip}

\section{ساختارِ یک پروژهٔ چندفایلی}

یک پروژهٔ متوسط معمولاً چنین سامان می‌یابد:

\begin{salamen}[language={}]
my-project/
    main.salam        // برنامهٔ اصلی، روالِ «ریشه» اینجاست
    math_utils.salam  // بستهٔ کمکیِ ریاضی
    text_utils.salam  // بستهٔ کمکیِ متن
\end{salamen}

\code{main.salam} با \code{واردسازی} ماژول‌های دیگر را می‌آورد و آن‌ها را به هم
می‌پیوندد. هنگامِ ساختِ برنامه، سلام به‌طورِ خودکار همهٔ فایل‌های واردشده را پیدا
و کامپایل می‌کند؛ کافی است فایلِ اصلی را به آن بدهید:

\begin{salamen}[language={}]
salam build main.salam --output=app.exe
\end{salamen}

\section{خلاصهٔ فصل}

\begin{itemize}
  \item \code{واردسازی} کتابخانه‌های آماده یا ماژول‌های خودمان را می‌آورد.
  \item \code{بسته} به یک فایل نام می‌دهد و آن را کتابخانه می‌کند.
  \item هر چیز به‌طورِ پیش‌فرض خصوصی است؛ \code{همگانی} آن را بیرون‌دیدنی
        می‌کند.
  \item با نامِ مستعار (\code{واردسازی fl "..."}) می‌توان ماژول را کوتاه صدا
        زد.
  \item با \code{@en} و \code{@fa} می‌توان به کدِ خود نام‌های دوزبانه داد.
  \item \code{lang()} زبانِ جاری را برمی‌گرداند.
\end{itemize}

\section{تمرین‌ها}

\exercise یک ماژول به نامِ \code{هندسه} بسازید که روال‌های همگانیِ \code{مساحت دایره}
و \code{محیط دایره} داشته باشد، و در فایلِ اصلی از آن استفاده کنید.

\exercise در ماژولِ هندسه یک روالِ کمکیِ \emph{خصوصی} بگذارید و نشان دهید که از
بیرون قابلِ دسترس نیست.

\exercise ماژولی بسازید که هم نامِ فارسی و هم انگلیسی برای روال‌هایش داشته باشد و
آن را یک بار از یک برنامهٔ فارسی و یک بار از یک برنامهٔ انگلیسی صدا بزنید.

\exercise ماژولِ \code{ریاضی}ِ استاندارد را وارد کنید و با \code{ریاضی.جذر} و
\code{ریاضی.توان} مساحتِ یک دایره و قطرِ آن را حساب کنید.

\exercise یک پروژهٔ سه‌فایلی طراحی کنید (اصلی + دو ماژول) و توضیح دهید هر فایل
چه مسئولیتی دارد.
